
\section{拟凝聚层的前推}

我们以讨论\trQCSheaf{}的\trPushforward{}与\trPullback、它们的性质以及一些应用，来为本章收尾。

设 $B\to A$ 是一个环映射。回想（由\mbox{练习~1.4.E}）：$\left(\cdot\otimes_B A,\ \cdot_B\right)$ 构成 $A$-\trModule{}的\trCategory{}与 $B$-\trModule{}的\trCategory{}之间的一个\trAdjointPair：我们有一个双射
\[
\Hom_A(N\otimes_B A,\,M)\cong\Hom_B(N,\,M_B)
\]
它在两个变元上都是\trFunctoriality{}的。这些构造对\trLocalization{}表现良好（在一种恰当的意义下），从而（常常）在概形上\trQCSheaf{}的\trCategory{}中行之有效（而且，事实上，它们在赋环空间上 $\OO$-\trModule{}的\trCategory{}中总是行之有效——见定义~14.5.6——尽管我们对此并不特别在意）。那个较为容易的构造（$M\mapsto M_B$）将变成我们的老朋友——\trPushforward。另一个（$N\mapsto A\otimes_B N$）则是\trPullback{}的一个亲戚，我不大情愿称它为「老朋友」。

让我们从\trPushforward{}开始，对它我们早已做过大半工作。

本节的主要寓意是：在「合理」的情势下，\trQCSheaf{}的\trPushforward{}仍是\trQC{}的，并且这一点可以借助上文定义的某个模构造来理解。我们从一个动机性的例子开始：

\textbf{14.4.A.（练习）}\quad 设 $\pi\colon\Spec A\to\Spec B$ 是仿射概形之间的一个\trMorphism，并设 $M$ 是一个 $A$-\trModule，于是 $\widetilde{M}$ 是 $\Spec A$ 上的一个\trQCSheaf。给出一个\trIsomorphism{} $\pi_*\widetilde{M}\tilde{\longrightarrow}M_B$。（提示：只有一种合理的做法：考察\trDistinguishedOpen。）

特别地，$\pi_*\widetilde{M}$ 是\trQC{}的。或许更重要的是，这蕴含：一个\trQCSheaf{}在\trAffine\trMorphism{}下的\trPushforward{}也是\trQC{}的。

\textbf{14.4.B.（练习）}\quad 若 $\pi\colon X\to Y$ 是一个\trAffine\trMorphism，证明 $\pi_*$ 是一个\trExact{}函子 $\mathit{QCoh}_X\to\mathit{QCoh}_Y$。

下面的结果推广了「\trQCSheaf{}在\trAffine\trMorphism{}下的\trPushforward{}也是\trQC{}的」这一事实。

\textbf{14.4.C.（必要练习：\mbox{练习~6.2.G} 之推论——「当 $\pi$ 为 QCQS 时，$\pi_*\colon\mathit{QCoh}_X\to\mathit{QCoh}_Y$」）}\footnote{译注：QCQS 为 quasicompact quasiseparated（拟紧拟分离）的缩写，系原书行话。}\quad 设 $\pi\colon X\to Y$ 是一个\trQuasicompact\trQuasiSeparated\trMorphism，$\scr{F}$ 是 $X$ 上的一个\trQCSheaf。证明 $\pi_*\scr{F}$ 是 $Y$ 上的一个\trQCSheaf。（若无\trQuasicompact{}或\trQuasiSeparated{}的假设便会出反例；见 [Stacks, tag 078C]。）

\trCoherentSheaf{}并不总能前推为\trCoherentSheaf。例如，考虑\trStructureMorphism{} $\pi\colon\mathbb{A}^1_k\to\Spec k$，它对应于 $k\to k[t]$。此时 $\pi_*\OO_{\mathbb{A}^1_k}$ 是对应于 $k[t]$ 的那个\trQCSheaf，而它不是一个\trFinitelyGenerated{}的 $k$-\trModule。但在好的情势下，\trCoherentSheaf{}确实会前推。例如：

\textbf{14.4.D.（练习）}\quad 设 $\pi\colon X\to Y$ 是局部\trNoetherian{}概形之间的一个\trFinite\trMorphism。若 $\scr{F}$ 是 $X$ 上的一个\trCoherentSheaf，证明 $\pi_*\scr{F}$ 是 $Y$ 上的一个\trCoherentSheaf。提示：先证明 $\pi_*\OO_X$ 是\trFiniteType{}的。

一旦我们定义了\trQCSheaf{}的\trCohomology，便会很快证明：若 $\scr{F}$ 是 $\mathbb{P}^n_k$ 上的一个\trCoherentSheaf，则 $\Gamma(\mathbb{P}^n_k,\scr{F})$ 是一个\trFiniteDimensional{}的 $k$-\trModule；而且，更一般地，若 $\scr{F}$ 是 $\Proj S_\bullet$ 上的一个\trCoherentSheaf，则 $\Gamma(\Proj S_\bullet,\scr{F})$ 是一个\trCoh{}的 $A$-\trModule（其中 $S_0=A$）。这是「由射影态射前推的\trCoherentSheaf{}也是\trCoherentSheaf」（定理~18.7.1(d)）这一事实的一个特殊情形。（「射影态射」这一观念——即 $\Proj S_\bullet\to\Spec A$ 的一个相对版本——将在 \S17.3 中定义。）

更一般地，在\trNoetherian{}假设下，由\trProper\trMorphism{}前推的\trCoherentSheaf{}也是\trCoherentSheaf（\trGrothendieck{}凝聚定理~18.9.1）。

\textbf{14.4.1.}\ \emph{\trhRemark{}。}\quad \trVectorBundle{}并不总能前推为\trVectorBundle，这是我们要在更大的（拟）凝聚层范畴中工作的另一个理由。（你能想出一个例子吗？提示：$\Spec k[t]/(t)\to\Spec k[t]$。）但在特别好的情形下它们可以前推，这正是\trCohomology{}与\trBaseChange{}定理~25.1.6 的目的之一。
